\chapter{Galois Cohomology}

\section{Group Cohomology}
Let $G$ be a discrete group, and let $p\colon EG \to BG$ be the universal
principal $G$-bundle over the classifying space $BG$. Its defining property
is that, for every paracompact Hausdorff space $X$, pulling back along maps
$X \to BG$ induces a bijection
\[
  [X, BG] \xrightarrow{\;\cong\;} \operatorname{Prin}_G(X)
\]
between homotopy classes of maps $X \to BG$ and isomorphism classes of
principal $G$-bundles over $X$.

In particular, let $\pi\colon P \to X$ be a principal $G$-bundle. Then there
is a map $f\colon X \to BG$, unique up to homotopy, and a $G$-equivariant map
$\tilde{f}\colon V \to EG$ such that the square
\[
\begin{tikzcd}
V \arrow[r, "\tilde{f}"] \arrow[d, "\pi"'] & EG \arrow[d, "p"] \\
X \arrow[r, "f"'] & BG
\end{tikzcd}
\]
is a pullback; that is, $V \cong f^{*}EG$ as principal $G$-bundles.

We note that since $G$ is discrete,
we have $\Pi_1(BG) \cong G$ and $\Pi_i(BG) = 0$ for $i > 1$.
So $BG$ is a $K(G, 1)$ space, and we have
\[
  H^i(G, M) \cong H^i(BG, M)
\]
for any $G$-module $M$.

Now, when we take a look at $H_{sing}^*(BG, \mathbb{Z})$,
it is much more complicated than $BG$
but it gives us a lot of information about $G$.
This gives a geometric motivation for group cohomology
as we study representations of $G$ to understand $G$ better.

We can study it when we want to understand
$G$, the Galois group of some field.
For example, when $G = \Z$,
we have $BG = S^1$ and $EG = \R$ and we can compute $H^*(S^1, \Z)$.
For another example, when $G = \Z/2\Z$,
we have $BG = \R P^{\infty}$ and $EG = S^{\infty}$ and we can compute
$H^*(\R P^{\infty}, \Z)$.
More generally, when $G = \Z/n\Z$,
we have $BG = L^{\infty}_n := S^{\infty}/(\Z/n\Z)$, the infinite lens space,
and $EG = S^{\infty}$. Here we regard $S^{\infty}$ as the unit sphere in
$\C^{\infty} = \bigcup_{k} \C^{k}$, and a generator of $\Z/n\Z$ acts by
multiplication by $\zeta = e^{2\pi i/n}$ on each coordinate. This action is
free, and $S^{\infty}$ is contractible, so $S^{\infty} \to L^{\infty}_n$ is a
universal principal $\Z/n\Z$-bundle. Using the CW structure on $L^{\infty}_n$
with one cell in each dimension, whose cellular chain complex is
\[
\Z \xleftarrow{\;0\;} \Z \xleftarrow{\;n\;} \Z \xleftarrow{\;0\;} \Z
\xleftarrow{\;n\;} \Z \xleftarrow{\;0\;} \cdots,
\]
we obtain
\[
H^{k}(L^{\infty}_n, \Z) \cong
\begin{cases}
\Z & k = 0, \\
\Z/n\Z & k > 0 \text{ even}, \\
0 & k \text{ odd}.
\end{cases}
\]
In fact, as a ring, $H^{*}(L^{\infty}_n, \Z) \cong \Z[x]/(nx)$ with
$\deg x = 2$. Setting $n = 2$ recovers $H^{*}(\R P^{\infty}, \Z)$, since
$L^{\infty}_2 = \R P^{\infty}$.

To study such objects,
we cannot find a good answer with topological methods, so we need to use algebraic methods.
We can study cohomology groups and Galois cohomology groups with algebraic methods.

What is a group cohomology?
For simplicity, we assume that $G$ is a finite group (although the profiniteness is enough).
Let $A$ be an abelian group with a $G$-action, i.e., a $G$-module.
Let $G\text{-}\Mod$ be the category of $G$-modules and $G$-equivariant homomorphisms.
We note that $G\text{-}\Mod$ is an abelian category with enough injectives.

Let's consider an example.
Let $\Z[G]$ be the group algebra of $G$ over $\Z$.
In fact, this is a canonical way of constructing a ring from a group.
Here are some properties of $\Z[G]$.
\begin{remark}\noindent
    \begin{enumerate}
        \item $\Z[G]$ is commutative if and only if $G$ is abelian.
        \item $\Z[G]$ is noetherian if and only if $G$ is finitely generated.
    \end{enumerate}
\end{remark}

\begin{remark}
    The algebra $\Z[G]$ is an augmented $\Z$-algebra
    with an augmentation $\epsilon\colon \Z[G] \to \Z$ map given by
    \[
        \sum_{g \in G} m_g g \longmapsto \sum_{g \in G} m_g.
    \]
    The augmentation ideal $I_G$ is defined as the kernel of $\epsilon$ so it is clear that it is a $2$-sided ideal of $\Z[G]$.
\end{remark}

We also note that $\Z[G]$ is a free $\Z$-module with a $\Z$-basis $G$.
So is $I_G$ with the following lemma.

\begin{lemma}
    The ideal $I_G$ has a $\Z$-basis $\{g - 1 \mid g \in G, g \neq 1\}$.
\end{lemma}
\begin{proof}
    Let $g \in G$ with $g \neq 1$. Then we have
    \[
        \epsilon(g - 1) = \epsilon(g) - \epsilon(1) = 1 - 1 = 0,
    \]
    so $g - 1 \in I_G$. Now, let $\sum_{g \in G} m_g g \in I_G$.
    Then we have
    \[
        0 = \epsilon\left(\sum_{g \in G} m_g g\right) = \sum_{g \in G} m_g.
    \]
    So we can write
    \[
        \sum_{g \in G} m_g g = m_1 + \sum_{g \neq 1} m_g g
        = -\sum_{g \neq 1} m_g + \sum_{g \neq 1} m_g g
        = \sum_{g \neq 1} m_g (g - 1).
    \]
    This shows that $\{g - 1\}_{g\neq 1}$ generates $I_G$ as a $\Z$-module.

    To show that it is linearly independent, suppose that
    $\sum_{g\neq 1} m_g (g - 1) = 0$. Then we have
    $\sum_{g\neq 1} m_g g + \sum_{g\neq 1} (-m_g)1 = 0$.
    Since the elements of $G$ are linearly independent in $\Z[G]$, we must have $m_g = 0$ for all $g\neq 1$.
\end{proof}

The action of $G$ is invariant on $I_G$ since $h\cdot (g-1) = hg - h = (hg - 1) - (h - 1) \in I_G$ for all $g, h \in G$.
So we have $\Z[G] = \Z \oplus I_G$ as $G$-modules.

As we mentioned, we can consider a ring $\Z[G]$ canonically from a group $G$.
Not surprisingly, we can consider a $G$-module $A$ as a $\Z[G]$-module
via $(\sum_{g \in G} m_g g) \cdot a = \sum_{g \in G} m_g (g \cdot a)$ for $a \in A$.
The converse is also true.
There is a composition of homomorphism $G \to \Z[G]^{\times} \to \mathrm{Aut}_{ab}(A)$
which gives a canonical $G$-action on $A$.

\begin{proposition}
    There is an equivalence of abelian categories between $G\text{-}\Mod$ and $\Z[G]\text{-}\Mod$.
\end{proposition}

Now we consider a functor $F: G\text{-}\Mod \to \mathsf{Ab}$
which sends $A$ to $A^G = \{a \in A \mid g\cdot a = a \text{ for all } g \in G\}$.
We may consider an abelian group $A$ as a $G$-module with a trivial action
so we have a functor $\mathrm{tr}: \mathsf{Ab} \to G\text{-}\Mod$.
\begin{claim}
We claim that $(\mathsf{tr}, F)$ is an adjoint pair of functors.
\end{claim}
\begin{proof}
    For any $A \in \mathsf{Ab}$ and $B \in G\text{-}\Mod$, we have a natural isomorphism
    \[
        \mathrm{Hom}_{G\text{-}\Mod}(\mathrm{tr}(A), B) \cong \mathrm{Hom}_{\mathsf{Ab}}(A, B^G).
    \]
    For any $G$-equivariant homomorphism $\phi: \mathrm{tr}(A) \to B$,
    we have a homomorphism $\psi: A \to B^G$ given by
    \[
        \psi(a) = \phi(a) \in B^G.
    \]
    Conversely, for any homomorphism $\psi: A \to B^G$,
    we have a $G$-equivariant homomorphism $\phi: \mathrm{tr}(A) \to B$ given by
    \[
        \phi(a) = \psi(a) \in B.
    \]
    It is clear that these two maps are inverses of each other.
\end{proof}

With the claim, we have a corollary.
\begin{corollary}
    The invariant functor $F: G\text{-}\Mod \to \mathsf{Ab}$ is left exact.
\end{corollary}

But note that $F$ is not right exact in general. For example, let $G = \Z/2\Z$
and take the augmentation sequence $0 \to I_G \to \Z[G] \to \Z \to 0$
which is exact with the definition of $I_G$.
But $0 \to I_G^G = 0 \to \Z[G]^G = \Z \to \Z^G = \Z \to 0$ is not exact.

With the left exactness of $F$, we can consider the right derived functor $R^iF$.
But we need to be careful; we need to have enough injectives in $G\text{-}\Mod$ to define the derived functor.